\chapter{Prompt Dictionary Concept Split (Phase 1.5)}
\label{chap:p1-5}

\section{Goal}

Florence-2 has no free-form VQA token, so every safety rule is answered by grounding a short noun
phrase. The pilot stored two phrases per rule in \texttt{RULE\_QUERIES} and used index~0 for the
answer and index~1 for the Day-9 ``rewording-robustness'' test, which reruns the rule with the second
phrasing and checks whether the answer or grounded box changes. The premise of that test is that
index~0 and index~1 are \emph{synonyms} --- two ways of saying the same thing --- so any change
measures sensitivity to \emph{wording}.

That premise is false for two of the three presence rules. For rule\_1 the two phrases are ``hard
hat'' and ``high-visibility vest''; for rule\_2 they are ``safety harness'' and ``fall-protection
lanyard.'' These are not synonyms --- they are \emph{distinct physical objects}. Only rule\_3's pair,
``guardrail'' and ``edge protection barrier,'' are genuine synonyms of one object. So the pilot's
rewording test, for rules 1 and 2, was silently measuring how the answer changes when you ask about a
\emph{different object}, and calling that ``robustness to rewording.'' The plan (item 3.5) prescribes
restructuring the query dictionary so each distinct object is its own concept with its own synonym
list, and so the rewording test only ever swaps validated synonyms within a concept.

\section{What Was Done}

\subsection{A concept/synonym structure}

A new \texttt{RULE\_CONCEPTS} maps each rule to \{concept: [synonyms]\}, disentangling the two things
the flat \texttt{RULE\_QUERIES} conflated:

\begin{center}
\begin{tabular}{p{2.0cm}p{5.5cm}p{5.0cm}}
\toprule
\textbf{rule} & \textbf{concepts (distinct objects)} & \textbf{synonyms per concept} \\
\midrule
rule\_1 & hard hat; high-visibility vest & \{hard hat, safety helmet\}; \{hi-vis vest, \ldots\} \\
rule\_2 & safety harness; fall-protection lanyard & \{safety harness, fall-arrest harness\}; \ldots \\
rule\_3 & guardrail (one concept) & \{guardrail, edge protection barrier\} \\
\bottomrule
\end{tabular}
\captionof{table}{The reviewed concept structure. rule\_1 and rule\_2 have two distinct concepts each;
rule\_3 has one concept with two genuine synonyms -- the only rule whose Day-9 ``rewording'' was
actually a rewording.}
\label{tab:p15-concepts}
\end{center}

A companion \texttt{RULE\_SATISFACTION} records whether a rule needs \emph{all} its concepts present or
\emph{any} one (rule\_1 ``all'': both hard hat and vest are expected PPE; rule\_2/rule\_3 ``any''),
with an explicit note that the pilot proxy approximated rule\_1 using only the hard-hat concept, so its
operative semantics were narrower than the recorded ``all'' intent. Helper functions expose the
structure (\texttt{concepts\_for}, \texttt{synonyms\_for}, \texttt{is\_synonym}), and
\texttt{reworded\_query} returns a validated same-concept synonym for the corrected rewording test ---
never a different concept.

\subsection{Frozen queries preserved}

\texttt{RULE\_QUERIES} is kept byte-for-byte (its index~0/index~1 phrases are unchanged), so every
existing script and the frozen Day-9 results reproduce exactly. The concept structure is additive: it
sits beside the frozen dictionary and is consumed only by code that opts into it. Ten unit tests,
written test-first, pin the disentanglement (rule\_3 synonyms group; rule\_1/rule\_2 objects
separate), the frozen-query preservation, and that \texttt{reworded\_query} only ever returns a
same-concept synonym.

\section{Results}

The structural claim --- that index~1 is a different \emph{object} for rules 1--2 --- is not just a
naming argument; it is measurable in where Florence-2 puts the box. Grounding the primary phrase versus
(a) a proposed same-concept synonym and (b) the frozen cross-concept phrase on 20 samples per rule, and
measuring agreement with the primary term's box (IoU; a missing or non-overlapping box scores 0):

\begin{center}
\begin{tabular}{p{2.0cm}p{4.6cm}p{3.0cm}p{2.8cm}}
\toprule
\textbf{rule} & \textbf{compared phrase} & \textbf{box agreement} & \textbf{grounded} \\
\midrule
rule\_1 & \emph{synonym} safety helmet          & \textbf{0.281} & 20/20 \\
rule\_1 & \emph{cross} high-visibility vest      & \textbf{0.000} & 18/20 \\
\midrule
rule\_2 & \emph{synonym} fall-arrest harness     & \textbf{0.426} & 13/20 \\
rule\_2 & \emph{cross} fall-protection lanyard   & \textbf{0.000} & 13/20 \\
\bottomrule
\end{tabular}
\captionof{table}{Grounding agreement with the primary concept phrase. The cross-concept phrases DO
ground a box (18/20, 13/20) --- but at a systematically different location: agreement is exactly 0.000.
True synonyms ground the same region (non-zero agreement).}
\label{tab:p15-iou}
\end{center}

The result is unusually clean. The cross-concept phrases are not failing to detect anything --- ``high-
visibility vest'' finds a box in 18 of 20 rule\_1 images, ``fall-protection lanyard'' in 13 of 20
rule\_2 images. But those boxes sit on a \emph{different object}, so their overlap with the hard-hat /
harness box is exactly zero. Swapping index~0 for index~1, as the pilot's rewording test did, therefore
relocated the grounding to a different object every time it grounded at all --- which is why any
answer or box change it produced cannot be attributed to wording. The true synonyms, by contrast, land
on the same object (agreement 0.281 and 0.426): imperfect box overlap, as expected from Florence-2's
phrasing-dependent box jitter, but unambiguously the \emph{same region}, not a different one.

\section{An Honest Note on Synonym Strength}

The same-concept agreements (0.281, 0.426) are modest, not near-unity, and that is worth stating
plainly rather than dressing up. Two effects are mixed in: Florence-2 genuinely returns somewhat
different box extents for different phrasings of the same object, and the proposed synonyms (``safety
helmet,'' ``fall-arrest harness'') are candidates, not yet blessed. The principled definition the
structure adopts is exactly this measurement: a proposed synonym is \emph{validated} when it grounds
substantially the same region as the primary phrase. By that definition ``safety helmet'' and ``fall-
arrest harness'' are weakly validated --- clearly the same object (non-zero, versus the cross term's
exact zero), but with enough box jitter that a scale run should either confirm them on more samples or
prefer tighter synonyms. The structure makes that a checkable property rather than an assumption; the
full study will run this check across the whole set before committing a synonym list.

\section{Standard vs. Non-Standard Choices}

\textbf{Standard.} Separating a controlled variable (wording) from a confound (object identity) is
basic experimental design; a robustness-to-rewording test that swaps objects is simply mis-specified.

\textbf{Non-standard, and the contribution.} Two choices. First, ``validated synonym'' is given an
operational, falsifiable definition --- same-region grounding, measured --- rather than left to
intuition, so the query dictionary can be audited empirically instead of argued about. Second, the
frozen queries and the entire pilot result set are preserved untouched beside the new structure, so the
mis-specified Day-9 rule\_1/rule\_2 numbers remain reproducible as exactly what they were, and the
corrected rewording test is a new, opt-in path rather than a silent redefinition of the old one.

\section{Implications}

The rewording-robustness test can now be run as a genuine phrasing-sensitivity measurement, swapping
only validated same-concept synonyms, with the corrected result reported separately from the frozen
(mis-specified) one. The concept structure also feeds two later needs: rule\_1's two-concept, ``all''-
satisfaction semantics is the hook for multi-object PPE checking (does a worker have both hard hat and
vest?), and the per-concept synonym lists give the native-VQA models a clean vocabulary to be prompted
with. The immediate carry-forward is a full-set synonym-validation pass to promote the weakly-validated
candidates or replace them.
